% Section 5 — Results
% Page budget: ~2 pages (the empirical payload)
%
% Subsections:
%   5.1 Headline: TCH vs every single-pipeline baseline
%       - Single table: per-pipeline Hit@1/Hit@5/MRR/PR-AUC, with TCH at the bottom
%       - Bootstrap CI delta column or footnote
%
%   5.2 Per-phase progression to the headline
%       - Cumulative effect of G1 (BiEncoder mixed-negs), G4 (agent Phase 3),
%         G7 (learned novelty)
%       - Single table showing v2f -> +G1 -> +G4 -> +G7
%
%   5.3 What the L1 stacker learned
%       - The HGB coefficient dominance (+8.221) figure or table
%       - Ablation: stacker without HGB collapses to PR-AUC=0.303
%
%   5.4 Three rerankers that failed (publishable negative pattern)
%       - G2 cross-encoder, Phase E agent rerank, G5 LLM judge
%       - Single table with each design + its delta vs cascade
%       - Cross-cutting interpretation: cascade's multi-retriever consensus is at a local optimum
%
%   5.5 Robustness
%       - G6 distractor sweep (Hit@5 -2% rel at 50% distractor ratio)
%       - G8 LOFO (F1 within 11% rel of in-distribution)
%       - Compact subsection — one paragraph each

\section{Results}
\label{sec:results}

\todo{The empirical payload. Write this first; everything else hangs off
      the numbers. Aim for ~1500 words spread across 4-5 subsections.
      Open with the headline table.}

\subsection{Headline metrics: TCH vs single-pipeline baselines}
\label{sec:results:headline}

\todo{Single per-pipeline table; one sentence each per row. Include the
      paired-bootstrap delta CIs. The table follows the
      technical-paper §7 Table 5 layout, compressed. End with: "TCH
      beats or ties every baseline on every metric, significantly so
      on Hit@5 and MRR; on Hit@1 the cascade ties BiEncoder within the
      95\% CI."}

% Headline table placeholder — fill from data/derived/global/.../comparison/v2g-final-models/final/tch_metrics.json
% \begin{table}
%   \centering
%   \caption{Per-pipeline metrics on the 1{,}008-window v2 in-distribution test split. Bold = best in column.}
%   \label{tab:headline}
%   \small
%   \begin{tabular}{lrrrr}
%   \toprule
%   Pipeline & Hit@1 & Hit@5 & MRR & PR-AUC \\
%   \midrule
%   \hgb{}                & ---     & ---     & ---     & \textbf{0.9998} \\
%   \bienc{}              & 0.695   & 0.789   & 0.729   & 0.287 \\
%   \hrrf{} rule          & 0.583   & 0.798   & 0.669   & 0.236 \\
%   \hrrf{} LLM           & 0.432   & 0.667   & 0.517   & 0.292 \\
%   \logseq{}             & 0.483   & 0.531   & 0.498   & 0.313 \\
%   \kgret{}              & 0.079   & 0.556   & 0.228   & ---     \\
%   \agent{}              & 0.386   & 0.436   & 0.405   & 0.243 \\
%   \midrule
%   \textbf{\tch{}}       & \textbf{0.722} & \textbf{0.912} & \textbf{0.794} & \textbf{0.9998} \\
%   \bottomrule
%   \end{tabular}
% \end{table}

\subsection{Where the lift came from: per-phase progression}
\label{sec:results:progression}

\todo{Single small table showing cumulative effect of G1 + G4 + G7.
      The narrative: G1 contributes a modest Hit@1 lift; G4 doubles
      novel recall; G7 turns that into a 4.9$\times$ novel-recall lift
      at preserved precision. Retrieval (Hit@5) stayed unchanged
      across the whole G-series, which is the headline finding.}

\subsection{Three rerankers that failed: a negative-result pattern}
\label{sec:results:rerankers}

\todo{This subsection earns the paper an extra novelty point. Three
      independent reranking designs (cross-encoder w/ BCE, agent rerank,
      LLM judge with structured JSON) all underperformed the cascade's
      overlap-rerank. The interpretation: multi-retriever consensus on
      top-1 dominates any single-LLM judgment on this dataset. Worth
      one paragraph + a 3-row table.}

\subsection{What the L1 stacker learned}
\label{sec:results:stacker}

\todo{One paragraph: HGB's coefficient dominates by ~30$\times$ the next
      largest; removing HGB collapses strict PR-AUC from 0.9998 to 0.303;
      the other features contribute on borderline-class windows where
      HGB alone is uncertain (+3.2 points inclusive PR-AUC).}

\subsection{Robustness}
\label{sec:results:robustness}

\todo{Two paragraphs:
        (a) Distractor sweep: Hit@5 drops only 2\% rel at 50\% distractor
            ratio; Hit@1 more sensitive (-14.2\% rel).
        (b) LOFO: novelty F1 within 11\% rel of in-distribution at the
            best threshold; precision robust across 19/27 families;
            recall splits by family structure.}
